% !PRINTED 113
\weberpage{113}
%<<CONT
and the number arising from it by the interchange of $i$ with $-i$,
and further by $c$ a rational number, we set
\[
\begin{aligned}
a + b i &= i^{\lambda}\, c^2 H_1^4\\
a - b i &= i^{-\lambda} c^2 H_2^4.
\end{aligned}
\]
From (8) it follows, by extraction of the $4$th root, whereby
a power of $i$ may enter as a factor,
\weq{15}{\begin{aligned}
(i, x_0) &= i^h \sqrt[4]{i}^{\lambda} \sqrt{c}\, H_1\\
(-i, x_0) &= i^{-h} \sqrt[4]{-i}^{\lambda} \sqrt{c}\, H_2\\
(-1, x_0) &= \sqrt{m},
\end{aligned}}
and from this, since $\sqrt[4]{i}$, $\sqrt{c}$ and $\sqrt{m}$ are cyclotomic
numbers, there results the correctness of our theorem in this case also.

Here too the further consideration of the structure of the expressions
occurring in (15) can serve to investigate more closely the composition
of the $x_0, x_1, x_2, x_3$ by the periods $\eta$; but this again can
lead to no other results than those already derived in \S.\,24.

We have thus herewith proved the theorem in general, \gesperrt{that
all Abelian fields of the third and fourth degrees
are cyclotomic fields, and that all these fields
can be represented rationally by the cyclotomic periods.}

\rule{3cm}{0.4pt}
